\documentclass{article}
\usepackage{graphicx} % Required for inserting images
\usepackage{hyperref} % For \href command
\usepackage{float}
\usepackage{amsmath}
\usepackage{listings}
\usepackage{xcolor}

\definecolor{codegreen}{rgb}{0,0.6,0}
\definecolor{codegray}{rgb}{0.5,0.5,0.5}
\definecolor{codepurple}{rgb}{0.58,0,0.82}
\definecolor{backcolour}{rgb}{0.95,0.95,0.92}

\lstdefinestyle{mystyle}{
    backgroundcolor=\color{backcolour},   
    commentstyle=\color{codegreen},
    keywordstyle=\color{magenta},
    numberstyle=\tiny\color{codegray},
    stringstyle=\color{codepurple},
    basicstyle=\ttfamily\footnotesize,
    breakatwhitespace=false,         
    breaklines=true,                 
    captionpos=b,                    
    keepspaces=true,                 
    numbers=left,                    
    numbersep=5pt,                  
    showspaces=false,                
    showstringspaces=false,
    showtabs=false,                  
    tabsize=2
}

\lstset{style=mystyle}

\title{Molecular Dynamics Simulation of a Lennard-Jones Liquid}
\author{Vedant Neema}
\date{\today}
  
\begin{document}
\maketitle

\section*{Description}
The program \cite{ljliquid} in this repository simulates a Lennard-Jones Liquid with variables:
$N = 100$ particles in a 2D square box with side, $L = 11$ ($\Rightarrow \rho = 0.83$), with velocities distributed as per $T = 1$

\section*{Dependencies}

I use the \href{https://www.raylib.com/}{\bf Raylib} Graphics Library v5.5-1 for the interface and \href{http://www.gnuplot.info/}{\bf gnuplot} v6.0 for the graphs.

\section*{How to run}
Look for the file \texttt{ljliquid.c}, get Raylib 5.0 or higher, then compile using the command
\begin{verbatim}
$> gcc ljliquid.c -lm -lraylib -o ljliquid
\end{verbatim}
This should produce a binary that can be run in the terminal like
\begin{verbatim}
$> ./ljliquid
\end{verbatim}

\newpage
\section*{Plots}
\begin{figure}[H]
    \centering
    \includegraphics[width=0.9\linewidth]{entempressure.png}
    \caption{Graph of energy, temperature and pressure}
    \label{fig:placeholder}
\end{figure}

\begin{figure}[H]
    \centering
    \includegraphics[width=0.9\linewidth]{rdf.png}
    \caption{Radial Distribution Function (rdf) sampled during the production run}
    \label{fig:placeholder}
\end{figure}

\newpage
\section*{Principles used}
Using the velocity Verlet algorithm, we get a order $\Delta t^4$ accuracy in position, $\Delta t^3$ in velocity and $\Delta t^2$ in Energy\cite{veve}:
\begin{itemize}
    \item Calculate $\mathbf {x} (t+\Delta t)=\mathbf {x} (t)+\mathbf {v} (t)\,\Delta t+\tfrac{1}{2}\,\mathbf {a} (t)\,\Delta t^{2}$.
    \item Derive $\displaystyle \mathbf {a} (t+\Delta t)$ from the interaction potential using $\displaystyle \mathbf {x} (t+\Delta t)$.
    \item Calculate $\displaystyle \mathbf {v} (t+\Delta t)=\mathbf {v} (t)+{\tfrac {1}{2}}\,{\bigl (}\mathbf {a} (t)+\mathbf {a} (t+\Delta t){\bigr )}\Delta t$.
\end{itemize}

\section*{Observations}
The rdf of the system looks like that of liquids. So, we can assume under the given conditions (T = 1, $\rho = 8.3$), any particles with the Lennard-Jones potential will be in the liquid state.

\section*{Code}
\lstinputlisting[language=C]{ljliquid.c}

\section*{Ending}
\begin{figure}[H]
    \centering
    \includegraphics[width=0.5\linewidth]{ending.png}
    \caption{Final Configuration Snapshot}
    \label{fig:placeholder}
\end{figure}

\begin{thebibliography}{9}
\bibitem{ljliquid}
    \href{https://github.com/Vedant36/ljliquid}{Repository containing ljliquid.c on Github}
\bibitem{veve}
    \href{https://en.wikipedia.org/wiki/Verlet_integration#Velocity_Verlet}{Velocity Verlet on Wikipedia}
\end{thebibliography}
\end{document}
